% Sumber angka: docs/ANALISIS_CHUNKING_T7.md dan results/eval_rag/ (16 Agustus 2026).
% Varian yang diadopsi V2 (keputusan bimbingan 18 Agustus 2026). Dasar adopsinya integritas
% korpus, bukan mutu penelusuran. Itu disengaja dan dinyatakan pada tabel maupun narasi.
\begin{table}[htbp]
\centering
\caption{Perbandingan empat strategi pemecahan dokumen atas korpus yang sama. ``Token terbuang'' adalah proporsi token korpus yang tidak pernah masuk ke vektor karena melewati jendela 256 token model \emph{embedding}. ``Akhir utuh'' adalah proporsi potongan yang berakhir pada kalimat lengkap.}
\label{tab:strategi-chunking}
\small
\begin{tabular}{@{}llrrrrr@{}}
\toprule
\textbf{Varian} & \textbf{Rancangan} & \textbf{Potongan} & \textbf{MRR} & \textbf{Hit@1} & \textbf{Token terbuang} & \textbf{Akhir utuh} \\
\midrule
V0 & 900 karakter, kontrol & 2.061 & 0,492 & 29,2\% & 8,23\% & 19,7\% \\
V1 & 900 karakter, batas kalimat & 2.248 & 0,468 & 29,2\% & 6,14\% & 83,5\% \\
\textbf{V2} & \textbf{500 karakter, batas kalimat} & \textbf{4.063} & 0,507 & 34,2\% & \textbf{0,00\%} & \textbf{85,3\%} \\
V3 & 256 token, batas kalimat & 2.587 & 0,150 & 3,3\% & \textbf{0,00\%} & 80,4\% \\
\bottomrule
\end{tabular}

\vspace{0.4em}
{\footnotesize \textbf{Jumlah potongan pada tabel ini dihitung sebelum dokumen kurasi manual dicabut}, sehingga tiap varian memuat 15 potongan yang kini tidak lagi diindeks; korpus produksi yang berlaku berisi \textbf{4.038} potongan. Angka itu dipertahankan apa adanya karena keempat varian diukur pada korpus yang sama, sehingga perbandingan antar-varian tetap sah; yang tidak sah adalah membandingkannya langsung dengan angka korpus produksi pada bab ini. \textbf{V2 adalah varian yang diadopsi ke produksi}. Dasar adopsinya \textbf{bukan} MRR-nya yang tertinggi melainkan \textbf{integritas korpus}: ia satu-satunya varian yang menuntaskan kedua cacat pemecahan sekaligus, yakni token terbuang 0,00\% dan akhir kalimat utuh 85,3\%. Varian V1 yang dipakai sebelumnya hanya menuntaskan cacat yang terlihat mata dan menyisakan 6,14\% token korpus yang tidak pernah dapat ditelusur karena ekornya tidak masuk ke vektor. Kolom MRR pada tabel ini \textbf{tidak dapat dipakai membandingkan antar-varian}, sebab label kebenaran metrik itu diturunkan dari teks potongan sehingga ikut bergerak ketika panjang potongan diubah; alasan yang sama membuat angka penelusuran pada korpus V1 dan V2 tidak sebanding satu sama lain. Uraiannya pada Subbab~\ref{sec:eval-chunking}.}
\end{table}
